% brando's dogfood of the BEAMER theme.
%
% //doc:fixture_doc compiles the article class and caught one half of 0.2.0's
% macro rename. It did not catch the other half, and fastverk's CI did:
%
%     deck.tex:15: Package xcolor Error: Undefined color `fvslate'.
%
% Two lessons, both narrow enough to be worth writing down.
%
% First, `fvslate` and `fvtertiary` were the SAME token under two names — the
% article class called it tertiary, the beamer theme called it slate — so
% unifying them to `brandmuted` broke a call site the article fixture never
% touches. A rule that emits three files needs a fixture per file, not one.
%
% Second, my first sweep grepped for `\fv[a-z]*` and so found the MACROS
% (\fvswatch, \fvheader) but not the COLOUR NAMES, which appear unprefixed inside
% braces: \color{fvslate}, text=fvfg. Grepping for the sigil misses half the
% surface; grep for the string.
\documentclass{beamer}
\usetheme{brandofixture}

\begin{document}

\begin{frame}{Colour}
  % Every colour the theme defines, used the way a real deck uses them: as a
  % colour NAME inside braces, which is the form the first sweep missed.
  {\color{brandbg}bg}\quad
  {\color{brandfg}fg}\quad
  {\color{brandaccent}accent}\quad
  {\color{brandmuted}muted}
\end{frame}

\begin{frame}{Structure}
  \begin{itemize}
    \item An itemize item, which the theme colours from \texttt{brandaccent}.
    \item A second, whose subitem marker uses \texttt{brandmuted}.
  \end{itemize}
  \begin{block}{A block}
    Block title and body, which is where the theme's \texttt{brandmuted!25} fill
    is exercised — a shade of a colour, so a renamed colour fails here too.
  \end{block}
\end{frame}

\end{document}
